\section{Discharging $\AND$: Hadamard as Ideal Membership}
\label{sec:hadamard}

This section proves the nonlinear part: the column Hadamard products that encode
$\AND$ (\Cref{sec:and}) and the carry Hadamards of the modular additions
(\Cref{sec:add}). The idea is to make the nonlinearity look like everything
else. Reading each cell's bits as the coefficients of a univariate over a small
subspace turns a bitwise $\AND$ into a polynomial \emph{ideal-membership} claim,
of exactly the kind the linear part already discharges. The Hadamard relations
then ride the same machinery --- one ideal check, the reading $\psia$, the
Step-4 sumcheck --- at the \emph{same} point, settled by the \emph{same} opening.
The only nonlinearity in SHA-256 thus costs one extra polynomial message and one
extra group in a sumcheck already running, and the commitment nothing beyond
reading a second functional off the one bound object.

\subsection{The obligation}
\label{sec:had-oblig}

Collect the $K_H = 14$ Hadamard relations --- the two $\AND$ products per round
and the carry Hadamard of each binary addition step --- as triples
\[
  \bp{u}^{(k)}\had\bp{v}^{(k)} = \bp{w}^{(k)},
  \qquad k = 1,\dots,K_H,
\]
where each operand is a fixed $\Ftwo$-linear combination of committed columns and
their shifts: for $\AND$, committed columns and their \emph{row}-shifts; for the
adders, the operands $\bp{a}+X\bp{c}$, $\bp{b}+X\bp{c}$, $\bp{c}+X\bp{c}$ of
\Cref{lem:binius}, in which $X\bp{c} = \SHL^{1}(\bp{c})$ is a \emph{within-cell}
bit-shift. The relation \eqref{eq:had-ideal} below is coefficient-wise among
operand bit-vectors, so the discharge is indifferent to the distinction; it
resurfaces only at binding (\Cref{rem:adder-trusted}). Writing
$u^{(k)}_b(x)\in\{0,1\}$ for bit $b$ of the operand at row $x$, the obligation is
\begin{equation}
  \label{eq:had-system}
  u^{(k)}_b(x)\,v^{(k)}_b(x) = w^{(k)}_b(x)
  \qquad \text{for all } k,\; b\in\{0,\dots,\wbit-1\},\; x\in\{0,1\}^{\nu}.
\end{equation}

\subsection{Bits as Lagrange coefficients: $\AND$ as ideal membership}
\label{sec:psi-z}

Identify the $\wbit$ bit positions with an $\Ftwo$-affine subspace $H\subset\K$
of dimension $\mu = \log_2\wbit$, with Lagrange basis $\{L_b\}$. Read a cell's
bit-vector $\bp{w}$ not as monomial coefficients (the reading $\psia$ uses) but
as the \emph{Lagrange} coefficients of
\[
  L(\bp{w})(X) = \sum_b w_b\,L_b(X),
  \qquad L(\bp{w})(h_b) = w_b \ \text{ for the base points } h_b\in H .
\]
Since $L(\bp{u}),L(\bp{v}),L(\bp{w})$ take the witness bits on the base domain
$H$, the per-bit products $u_b v_b = w_b$ of \eqref{eq:had-system} hold for all
$b$ exactly when $L(\bp{u})\,L(\bp{v}) - L(\bp{w})$ vanishes on all of $H$ ---
exactly when it lies in the ideal generated by $Z_H(X) = \prod_{h\in H}(X-h)$:
\begin{equation}
  \label{eq:had-ideal}
  \bp{u}^{(k)}\had\bp{v}^{(k)} = \bp{w}^{(k)}
  \quad\Longleftrightarrow\quad
  L(\bp{u}^{(k)})\,L(\bp{v}^{(k)}) - L(\bp{w}^{(k)}) \;\in\; \langle Z_H\rangle .
\end{equation}
This is the shape of the linear part's obligations (membership in $\langle 0\rangle$
or $\langle X\rangle$), with the extra generator $Z_H$. Batching the $K_H$
relations by a challenge $\gamma_{\mathrm{rel}}$ and the rows by $\eq(\cdot;r)$,
the whole nonlinear obligation is the single membership claim
\begin{equation}
  \label{eq:had-zerocheck}
  \sum_{x\in\{0,1\}^{\nu}} \eq(x;r)
  \sum_{k} \gamma_{\mathrm{rel}}^{\,k}
  \Bigl(L(\bp{u}^{(k)})\,L(\bp{v}^{(k)}) - L(\bp{w}^{(k)})\Bigr)(X,x)
  \;\in\; \langle Z_H\rangle ,
\end{equation}
where $L(\bp{u}^{(k)})(X,x) = \sum_b u^{(k)}_b(x)L_b(X)$.

\subsection{Discharging it on the main pipeline}
\label{sec:had-pipeline}

Membership \eqref{eq:had-zerocheck} is checked by the same three steps as the
linear ideal check, and at the same points.

\paragraph{The ideal-check message.} The prover sends the witness for membership
in $\langle Z_H\rangle$: the off-subspace values of
\[
  R_0(X) = \sum_{x}\eq(x;r)\sum_k\gamma_{\mathrm{rel}}^{\,k}
  \bigl(L(\bp{u}^{(k)})\,L(\bp{v}^{(k)}) - L(\bp{w}^{(k)})\bigr)(X,x),
  \qquad \deg R_0 \le 2(\wbit-1).
\]
Two economies make this cheap. On the \emph{base} domain $X\in H$ the operands
take the $\{0,1\}$ witness bits and $R_0$ \emph{vanishes for honest data} by
\eqref{eq:had-ideal} --- so the prover transmits only the $\wbit-1$
\emph{off-subspace} values, the genuine content of the $\langle Z_H\rangle$
witness. And the dominant accumulation over the $2^{\nu}$ rows, restricted to the
base domain, is bit arithmetic; only the off-subspace lift touches $\K$. The
message is absorbed into the transcript before $\alpha$.

\paragraph{The reading at $\alpha$.} The same point $\alpha$ that reads the
linear constraints settles membership: a degree-$<\wbit$ polynomial lies in
$\langle Z_H\rangle$ iff it vanishes on $H$, and the verifier tests this at
$X = \alpha$ via the shipped off-subspace values. Reading the operands at
$\alpha$ is the Lagrange instance of \Cref{def:projection},
\[
  \psiz(\bp{w}) = L(\bp{w})(\alpha) = \sum_b L_b(\alpha)\,w_b \in\K ,
\]
the weighting $\xi_b = L_b(\alpha)$ --- the analogue of $\psia$'s
$\xi_b = \alpha^b$, read at the \emph{same} $\alpha$. The product term then joins
the Step-4 sumcheck (\Cref{sec:sumcheck}) as one degree-$3$ group,
\[
  \sum_{x}\eq(x;r)\sum_k\gamma_{\mathrm{rel}}^{\,k}
  \bigl(\psiz(\bp{u}^{(k)})(x)\,\psiz(\bp{v}^{(k)})(x) - \psiz(\bp{w}^{(k)})(x)\bigr),
\]
whose claimed sum is $R_0(\alpha)$, reconstructed from the off-subspace message.
The bit dimension never becomes a sumcheck variable: the $\mu$ bit positions are
folded once, at $\alpha$, by the Lagrange reading.

\paragraph{The shared sumcheck.} Riding the Step-4 sumcheck, the group is reduced
over the same $\nu$ row variables to the same point $r^*$, yielding the closing
operand evaluations
\begin{equation}
  \label{eq:operand-evals}
  \psiz\bigl(\bp{u}^{(k)}\bigr)(r^*),\quad
  \psiz\bigl(\bp{v}^{(k)}\bigr)(r^*),\quad
  \psiz\bigl(\bp{w}^{(k)}\bigr)(r^*),
\end{equation}
whose combination
$\sum_k\eq(k;\gamma_{\mathrm{rel}})\,(a^{(k)}b^{(k)}-c^{(k)})$ the verifier checks
against the group's reduced claim. There is no second sumcheck and no second
point.

\begin{lemma}[The Hadamard reduction is sound]
\label{lem:psi-z-sound}
If \eqref{eq:had-system} holds then \eqref{eq:had-zerocheck} holds and $R_0$
vanishes on $H$, so the check passes for every $\gamma_{\mathrm{rel}},r,\alpha$.
Conversely, if some equation of \eqref{eq:had-system} fails, then over uniform
$\gamma_{\mathrm{rel}}, r, \alpha$ the check passes with probability at most
$\frac{2(\wbit-1)}{|\K|}+\frac{K_H-1}{|\K|}+\frac{\nu}{|\K|}$.
\end{lemma}

\begin{proof}[Idea]
A failed equation makes
$\bp{u}^{(k)}\had\bp{v}^{(k)}-\bp{w}^{(k)}$ nonzero at some base point, so
$L(\bp{u}^{(k)})L(\bp{v}^{(k)})-L(\bp{w}^{(k)})$ does not vanish on $H$ and $R_0$
is a nonzero univariate of degree $\le 2(\wbit-1)$; it is nonzero at the random
$\alpha$ except with probability $\frac{2(\wbit-1)}{|\K|}$ (Schwartz--Zippel).
The relation batch and the $\eq(\cdot;r)$ row check add the remaining terms by
the standard arguments. The dominant term is the analogue of the linear ideal
check's $\frac{\wbit-1}{|\K|}$, since $\psia$ and $\psiz$ are the same evaluation
at $\alpha$, read in two bases.
\end{proof}

The Lagrange basis is needed precisely because of \Cref{rem:and-not-survive}: it
keeps the bit-products $u_b v_b$ separated by position, so the verifier's closing
check is a genuine product of per-position bits, where $\psia$ of
$\bp{u}\had\bp{v}$ would scramble them. Both readings live at the one point
$\alpha$ and are read off one committed object, as the binding now shows.

\begin{remark}[The $\mathrm{GF}(2^8)$ acceleration]
\label{rem:gf8}
The base-domain accumulation can run the additive NTT that evaluates $R_0$ over
an \emph{embedded} $\mathrm{GF}(2^8)$ subspace with a byte-lookup table, rather
than over the monomial subspace of $\K$. This is a pure performance variant: the
embedding $\mathrm{GF}(2^8)\hookrightarrow\K$ is a field homomorphism, so the
message computed in $\mathrm{GF}(2^8)$ equals its image in $\K$, and
\Cref{lem:psi-z-sound} is unchanged. The deployed prover uses it.
\end{remark}

\subsection{Binding the discharge to the commitment}
\label{sec:had-binding}

The reduction has left the operand evaluations \eqref{eq:operand-evals} at $r^*$.
Two facts let these ride the \emph{same} opening that settles the linear part.

\emph{$\AND$ operands are $\Ftwo$-linear in column row-shifts.} Each $\AND$
operand is an $\XOR$ of committed columns and their row-shifts, and $\psiz$ is
$\Ftwo$-linear, so $\psiz(\bp{u}^{(k)})(r^*)$ is the $\Ftwo$-sum of the
per-column \emph{pair-evals} $\psiz(\col{c}^{\downarrow\Delta})(r^*)$ --- exactly
as the linear path needs the $\psia$ pair-evals, and through the same row-shift
machinery.

\emph{Pair-evals fold into the multipoint eval.} The
$\psiz(\col{c}^{\downarrow\Delta})(r^*)$ are appended to the main multipoint
evaluation (\Cref{sec:multipoint}) as pointed-shift claims and reduced ---
alongside the $\psia$ claims --- to the single point $r_0$. There the opening
binds, for each column, the entire bit-slice fold $a'_g$
(\Cref{sec:commitment}). Because one bound object carries all bit slices and both
readings live at the same $\alpha$, the verifier reads \emph{both} weightings off
it:
\[
  \psia(a'_g) = \mle{\psia(w_g)}(r_0)\ \ (\text{linear}),
  \qquad
  \psiz(a'_g) = \mle{\psiz(w_g)}(r_0)\ \ (\text{discharge}).
\]
Concretely, the open exposes its column-batching challenge $\gamma$, and the
single binding check
\begin{equation}
  \label{eq:psi-z-binding}
  \psiz(a') \;\stackrel{?}{=}\; \sum_{g} \gamma_g\cdot\mle{\psiz(w_g)}(r_0),
  \qquad a' = \sum_g\gamma_g a'_g,
\end{equation}
ties the discharge's evaluations to the committed bit slices. A final \emph{tie}
recombines the bound pair-evals into the operand evaluations
\eqref{eq:operand-evals} and checks them against the group's closing claim. This
is the entire cost of binding: one extra weight-vector reading
\eqref{eq:psi-z-binding} of the already-bound $a'$, no new opening and no new
point.

\begin{remark}[Adder operands are trusted]
\label{rem:adder-trusted}
The $\AND$ operands bind as above. The adder operands look more delicate only
because of $X\bp{c} = \SHL^{1}(\bp{c})$, a within-cell shift of the carry. Two of
the three are not even shifts: the linear identity \eqref{eq:add-linear} forces
$X\bp{c} = \bp{D}$ (with $\bp{D} = \bp{s}+\bp{a}+\bp{b}$), so
$\bp{a}+X\bp{c} = \bp{s}+\bp{b}$ and $\bp{b}+X\bp{c} = \bp{s}+\bp{a}$ are plain
$\XOR$s of committed columns. Only $\bp{c}+X\bp{c} = \SHR^{1}(\bp{D})+\bp{D}$,
with the carry-out bit $\beta$, keeps a within-cell shift --- which binds the way
a row-shift does, one dimension over, by reading the bound fold with
\emph{shifted} Lagrange weights $\psiz(\SHL^{1}\bp{c}) = \sum_b L_{b+1}(\alpha)a'_b$.
The deployed prover does not yet perform this bit-shift read-off and does not
commit $\beta$, so it ships the adder $\psiz$ parents \emph{trusted}
(honest-prover-supplied) and folds them into the tie. Committing $\beta$ and
adding the bit-shift binding closes the gap; it is localized to the adder
operands and leaves the $\AND$ relations and the entire linear part
unconditional. A recorded follow-up.
\end{remark}

\subsection{Summary}
\label{sec:had-summary}

The Hadamard part of SHA-256 --- its only nonlinearity --- is an
ideal-membership claim \eqref{eq:had-ideal} in $\langle Z_H\rangle$, obtained by
reading cell bits as Lagrange coefficients. It is discharged by the linear
part's own machinery: one ideal-check message (the off-subspace values of
$R_0$), one extra degree-$3$ group in the Step-4 sumcheck at the shared $\alpha$,
reduced to the shared $r^*$, with soundness $O((\wbit + K_H + \nu)/|\K|)$
(\Cref{lem:psi-z-sound}). It binds by the single reading
\eqref{eq:psi-z-binding} on the shared bit-slice opening. The discharge thus
costs, at the commitment layer, nothing beyond reading the cell in a second
basis off the one object the opening already binds --- the property the next
section's commitment provides.
